% =============================================================================
\chapter{Arquitetura em três camadas}
\label{cap:arquitetura}
% =============================================================================

\section{As camadas e o que cabe a cada uma}

O Morphe se divide em três camadas, uma por máquina ou por tipo de circuito
(Figura~\ref{fig:camadas}). Cada uma é escrita numa linguagem e tem uma
responsabilidade que as outras não repetem.

\begin{figure}[htbp]
  \centering
  % As tres camadas e os modulos de cada uma.
\begin{tikzpicture}[
    mod/.style={caixa, minimum width=27mm, minimum height=9mm, font=\scriptsize},
    faixa/.style={draw, rounded corners=4pt, inner sep=4mm},
  ]
  % --- Cliente -----------------------------------------------------------
  \node[mod, cliente] (jan)  at (0,0)     {janelas\\\texttt{*\_window.py}};
  \node[mod, cliente] (blo)  at (30mm,0)  {\texttt{blocos.py}\\sinais longos};
  \node[mod, cliente] (dsp)  at (60mm,0)  {\texttt{dsp\_core.py}\\ponto fixo e referência};
  \node[mod, cliente] (proj) at (90mm,0)  {projetistas\\\texttt{fir/iir\_design}};
  \node[mod, cliente] (pro)  at (120mm,0) {\texttt{morphe\_protocol}\\cliente TCP};

  % --- Servidor ----------------------------------------------------------
  \node[mod, hps, minimum width=52mm] (srv) at (60mm,-24mm)
        {\texttt{morphe\_server.c}\\valida, copia, dispara, espera, lê};
  \node[mod, hps] (mk) at (0,-24mm) {marca\\\texttt{fpga-preparada}};
  \draw[seta] (srv.west) -- node[rotulo, above]{consulta} (mk.east);

  % --- FPGA --------------------------------------------------------------
  \node[mod, fpga] (conv) at (15mm,-48mm)  {\texttt{conv1d}\\convolução};
  \node[mod, fpga] (fir)  at (45mm,-48mm)  {\texttt{conv1d}\\FIR};
  \node[mod, fpga] (fft)  at (75mm,-48mm)  {IP FFT II\\FFT e IFFT};
  \node[mod, fpga] (iir)  at (105mm,-48mm) {\texttt{iir\_cascade}\\IIR};

  % Faixas, todas da mesma largura
  \begin{scope}[on background layer]
    \node[faixa, draw=cCliente, fill=cCliente!4, fit=(jan)(pro),
          label={[anchor=south west, inner sep=1pt, font=\scriptsize\bfseries, text=cCliente]north west:{PC --- Python}}] {};
    \node[faixa, draw=cHps, fill=cHps!4, fit=(mk)(srv)(pro.east |- srv),
          label={[anchor=south west, inner sep=1pt, font=\scriptsize\bfseries, text=cHps]north west:{HPS --- C, Linux}}] {};
    \node[faixa, draw=cFpga, fill=cFpga!4, fit=(conv)(iir)(jan.west |- conv)(pro.east |- iir),
          label={[anchor=south west, inner sep=1pt, font=\scriptsize\bfseries, text=cFpga]north west:{FPGA --- Verilog}}] {};
  \end{scope}

  \draw[seta] (jan) -- (blo);
  \draw[seta] (blo) -- (dsp);
  \draw[seta] (proj) -- (dsp);
  \draw[seta] (dsp.north) to[bend left=25] (pro.north);
  \draw[<->, thick, cRede] (pro.south) |- node[rotulo, pos=0.3, left]{TCP\\5000} (srv.east);
  \foreach \m in {conv, fir, fft, iir}
    \draw[seta] (srv.south) -- ++(0,-5mm) -| (\m.north);
  \node[rotulo, right] at ($(srv.south)+(1mm,-2.5mm)$) {\texttt{/dev/mem}};
\end{tikzpicture}

  \caption{As três camadas e os módulos de cada uma.}
  \label{fig:camadas}
\end{figure}

\begin{table}[htbp]
  \centering
  \small
  \caption{Responsabilidades de cada camada.}
  \label{tab:camadas}
  \begin{tabularx}{\textwidth}{@{}lXX@{}}
    \toprule
    & \textbf{Faz} & \textbf{Não faz} \\
    \midrule
    \textbf{Cliente} \newline \footnotesize Python, no PC
      & gera ou lê o sinal; projeta filtros; converte para ponto fixo; divide
        sinais longos em blocos; calcula a referência em ponto flutuante;
        compara e mostra
      & não sabe endereços de memória nem detalhes da FPGA \\
    \addlinespace
    \textbf{Servidor} \newline \footnotesize C, no HPS da placa
      & valida a requisição; copia os dados para as memórias da FPGA; dispara o
        acelerador; espera; devolve o resultado
      & não converte formatos numéricos (com uma exceção na FFT, seção
        \ref{sec:acel-fft}); não guarda nada entre requisições \\
    \addlinespace
    \textbf{Aceleradores} \newline \footnotesize Verilog, na FPGA
      & fazem a conta, em ponto fixo, sobre blocos de tamanho fixo
      & não conhecem a rede nem o tamanho do sinal que o aluno pediu \\
    \bottomrule
  \end{tabularx}
\end{table}

\subsection*{Três princípios que atravessam as camadas}

\paragraph{O servidor não guarda estado.} Cada operação chega completa --- os
dados e todos os parâmetros --- e sai completa. Não há sessão, reserva de placa
nem memória da requisição anterior. É isso que permite que vários clientes usem
a mesma placa intercalando operações, e que um cliente possa mudar de placa de
uma operação para a outra sem nada a transferir.

\paragraph{O ponto fixo é decidido no cliente.} O cliente converte o sinal para
o formato do acelerador (Q15.16 ou Q15.8) antes de enviar, e converte o
resultado de volta. O servidor copia inteiros de 32 bits sem interpretá-los,
com uma única exceção: aplicar o expoente de bloco da FFT, que só ele conhece
(seção~\ref{sec:acel-fft}). Fora isso, toda decisão numérica --- normalizar a
entrada, onde saturar, que fator aplicar à IFFT --- mora num lugar só, em
Python, onde é fácil de ler e de testar.

\paragraph{O hardware tem tamanho fixo; o software se adapta.} Os aceleradores
foram sintetizados para 1024 amostras. O servidor completa com zeros o que o
cliente não mandou e devolve só as amostras que fazem sentido; o cliente, por
sua vez, divide sinais maiores em blocos de 1024. Nenhum dos dois lados exige
que o aluno pense em 1024.

\subsection*{Os números de configuração existem em dois lugares}

Os limites do hardware --- 1024 amostras, 16 seções de IIR, os bits
fracionários de cada formato --- estão escritos em \arq{C/morphe_config.h}, para
o servidor, e repetidos em \arq{Python/morphe_config.py}, para o cliente. Os
dois arquivos precisam mudar juntos. Os endereços das memórias e registradores,
ao contrário, não são escritos à mão: vêm do \arq{C/hps_0.h}, gerado a partir do
projeto do hardware (seção~\ref{sec:pontes}).

\section{O caminho de uma operação}
\label{sec:caminho}

A Figura~\ref{fig:sequencia} acompanha uma convolução do clique na janela até o
gráfico. As outras operações seguem o mesmo roteiro, mudando só as memórias e o
formato numérico.

\begin{figure}[htbp]
  \centering
  % Diagrama de sequencia de uma operacao (ex.: convolucao).
\begin{tikzpicture}[
    ator/.style={caixa, minimum width=26mm, font=\scriptsize},
    msg/.style={->, thick},
    nota/.style={font=\scriptsize, align=left, fill=white, inner sep=1pt},
  ]
  \def\xa{0}\def\xb{36mm}\def\xc{74mm}\def\xd{112mm}
  \def\fim{-112mm}

  \node[ator, cliente] (a) at (\xa,0) {Janela\\(\texttt{conv\_window})};
  \node[ator, cliente] (b) at (\xb,0) {\texttt{morphe\_protocol}};
  \node[ator, hps]     (c) at (\xc,0) {\texttt{morphe\_server}};
  \node[ator, fpga]    (d) at (\xd,0) {\texttt{conv1d}\\e memórias};

  \foreach \n in {a,b,c,d}
    \draw[densely dashed, black!40] (\n.south) -- (\n.south |- 0,\fim);

  % Atividade no servidor
  \fill[cHps!25] ($(\xc,-26mm)+(-1mm,0)$) rectangle ($(\xc,-100mm)+(1mm,0)$);

  \draw[msg] (\xa,-12mm) -- node[above, rotulo]{\(x[n],\ h[n]\) em ponto flutuante} (\xb,-12mm);
  \draw[msg] (\xb,-17mm) arc[start angle=90, end angle=-90, radius=2.5mm]
             node[right, nota, pos=0.5, xshift=1mm] {codifica em Q15.16};
  \draw[msg, cRede] (\xb,-26mm) -- node[above, rotulo]{\texttt{connect}; cabeçalho de 20\,B + \emph{payload}} (\xc,-26mm);
  \draw[msg] (\xc,-32mm) arc[start angle=90, end angle=-90, radius=2.5mm]
             node[right, nota, pos=0.5, xshift=1mm] {valida magic, versão,\\opcode, tamanhos, marca};
  \draw[msg] (\xc,-44mm) -- node[above, rotulo]{escreve \(x\) e \(h\) nas SRAMs} (\xd,-44mm);
  \draw[msg] (\xc,-52mm) -- node[above, rotulo]{\emph{start} = 1 (ponte leve)} (\xd,-52mm);
  \draw[msg, dashed] (\xd,-60mm) -- node[above, rotulo]{lê \emph{done} até 1 (limite 5\,s)} (\xc,-60mm);
  \node[nota, right] at (\xd,-57mm) {calcula\\\(\approx\)214\,ms};
  \draw[msg, dashed] (\xd,-70mm) -- node[above, rotulo]{lê \(y\) da SRAM de saída} (\xc,-70mm);
  \draw[msg, cRede] (\xc,-82mm) -- node[above, rotulo]{cabeçalho de resposta + \(y\); \texttt{close}} (\xb,-82mm);
  \draw[msg] (\xb,-88mm) arc[start angle=90, end angle=-90, radius=2.5mm]
             node[right, nota, pos=0.5, xshift=1mm] {decodifica, remove\\o \emph{padding}};
  \draw[msg] (\xb,-100mm) -- node[above, rotulo]{\(y\) da placa \(\times\) \(y\) de referência} (\xa,-100mm);
\end{tikzpicture}

  \caption{O caminho de uma operação, do clique na janela ao gráfico.}
  \label{fig:sequencia}
\end{figure}

\begin{enumerate}
  \item \textbf{A janela entrega os sinais.} \(x[n]\) e \(h[n]\) chegam ao
    módulo de protocolo como vetores em ponto flutuante, do tamanho que o aluno
    escolheu.

  \item \textbf{O cliente codifica.} Cada amostra vira um inteiro de 32 bits em
    Q15.16: \(q = \operatorname{round}(v \cdot 2^{16})\), saturado na faixa do
    formato. Se o sinal passa de 1024 amostras, a divisão em blocos acontece
    antes deste passo, e os passos seguintes se repetem por bloco.

  \item \textbf{Uma conexão por operação.} O cliente abre uma conexão TCP com a
    placa, envia um cabeçalho de 20 bytes (operação, tipo dos dados, tamanhos) e
    os dados logo em seguida (capítulo~\ref{cap:protocolo}).

  \item \textbf{O servidor valida.} Confere a assinatura e a versão do
    protocolo, a operação, os tamanhos contra os limites do hardware e se a FPGA
    foi preparada pelo \arq{morphe-up.sh} (capítulo~\ref{cap:servidor}). Qualquer
    falha vira uma resposta de erro com uma mensagem legível, e a FPGA não é
    tocada.

  \item \textbf{O servidor escreve nas memórias.} Copia \(x\) e \(h\) para as
    memórias de entrada do acelerador, através da ponte HPS\(\to\)FPGA,
    completando com zeros até 1024 amostras.

  \item \textbf{O servidor dispara.} Escreve 1 no registrador \texttt{start}
    do acelerador, através da ponte leve. O acelerador detecta a subida desse
    sinal e começa.

  \item \textbf{O acelerador calcula} e, ao terminar, leva o sinal
    \texttt{done} a 1. A convolução de 1024 por 1024 amostras leva cerca de
    214\,ms; a FFT de 1024 pontos, cerca de 21\,ms (tempos de ida e volta
    medidos pelo \arq{morphe_ping} em 15/09/2026, com a placa sem outros
    clientes).

  \item \textbf{O servidor espera.} Lê o \texttt{done} a cada 50\,\textmu s. Se ele
    não subir em 5\,s, a operação é abandonada com o erro de tempo esgotado ---
    o sintoma típico de um bitstream que parou de funcionar
    (seção~\ref{sec:aceleradores}).

  \item \textbf{O servidor lê e responde.} Copia as \(n_x + n_h - 1\) amostras
    úteis da memória de saída, devolve o \texttt{start} a 0 e envia o cabeçalho
    de resposta e os dados. Fecha a conexão.

  \item \textbf{O cliente decodifica e compara.} Converte de volta para ponto
    flutuante, descarta o que for enchimento de zeros e calcula a referência em
    ponto flutuante com o NumPy. A janela mostra as duas e o erro entre elas.
\end{enumerate}

Dois fatos decorrem desse roteiro e aparecem na prática.

O tempo de uma operação \textbf{não depende do tamanho do sinal pedido}: o
acelerador sempre processa os blocos completos de 1024 amostras, e os zeros de
enchimento custam o mesmo que amostras de verdade. Pela construção do
\texttt{conv1d}, uma convolução de 10 por 10 amostras custa o mesmo que uma de
1024 por 1024.

Como o servidor atende \textbf{uma conexão de cada vez}, duas pessoas na mesma
placa não recebem erro: a segunda espera a primeira terminar, e o tempo de cada
uma soma. O capítulo~\ref{cap:placas} trata disso.

\section{Onde está cada camada no repositório}

\begin{center}
\small
\begin{tabular}{@{}ll@{}}
  \toprule
  \textbf{Camada} & \textbf{Onde} \\
  \midrule
  Cliente      & \arq{Python/} (aplicativo e módulos); testes em \arq{Python/ferramentas/} \\
  Servidor     & \arq{C/} (\arq{morphe_server.c} e cabeçalhos) \\
  Aceleradores & \arq{Quartus/} (RTL, projeto do Platform Designer, bitstream) \\
  Operação     & \arq{morphe-up.sh}, \arq{instala-quartus.sh}, \arq{deploy.sh}, na raiz \\
  \bottomrule
\end{tabular}
\end{center}

\noindent A árvore completa está no apêndice~\ref{ap:repositorio}.
